% =============================================================
%  SECTION: Module 4 — Performance and Optimization
%  ASSIGNED TO: [Person Name]
%
%  GUIDING QUESTIONS TO ANSWER:
%   WHAT  — What is the O(N²) bottleneck? What is a Verlet neighbor list?
%           What is NumPy vectorization?
%   WHY   — Why does the naive force loop become a problem at N=500?
%           Why does the skin radius allow infrequent rebuilds?
%           Why does the shifted force potential improve energy conservation?
%   HOW   — How is the neighbor list built? When is it rebuilt?
%           How does vectorization replace Python for-loops?
%  EXTRAS — Benchmark: show timing comparison naive vs optimized (optional).
%           Explain the shifted-force potential and how it differs from
%           the shifted potential in Module 1.
% =============================================================

\newpage
\section{Module 4: Performance and Optimization}

\subsection{What Are We Doing?}
% Refactoring the force calculation to (1) skip distant pairs using a neighbor
% list, and (2) use NumPy array operations instead of Python loops.
% [WRITE HERE]

\subsection{Why Are We Doing It?}

\subsubsection{The \texorpdfstring{$O(N^2)$}{O(N²)} Bottleneck}
% With N atoms, naive force loop: N*(N-1)/2 pair checks per step.
% For N=500: ~125,000 checks every timestep. Most are beyond r_c and contribute zero.
% Show or cite the scaling: double N → 4x slower.
% [WRITE HERE]

\subsubsection{Why Does the Skin Radius Work?}
% Atoms move slowly relative to r_skin - r_c.
% Only rebuild when max displacement > (r_skin - r_c)/2.
% Analogy: you don't update your contact list every second — only when someone moves.
% [WRITE HERE — good place for physical intuition]

\subsubsection{Why the Shifted-Force Potential?}
% Even with a shifted potential, the force magnitude has a jump at r_c.
% This jump acts as a source/sink of energy every time a pair crosses the boundary.
% The shifted-force potential sets both phi and F smoothly to zero at r_c.
% [WRITE HERE]

\subsection{How Is It Implemented?}

\subsubsection{Verlet Neighbor List}

All pairs $(i, j)$ with separation $r_{ij} < r_s$ are stored, where the skin radius is:
\begin{equation}
    r_s = r_c + \delta_{\text{skin}}
    \label{eq:rskin}
\end{equation}
We used $r_c^* = 2.5$ and $r_s^* = 3.2$, giving a skin thickness of $0.7\,\sigma$.

The list is rebuilt only when any atom has displaced more than half the skin thickness from its last recorded position:
\begin{equation}
    \Delta r_{\max} > \frac{r_s - r_c}{2}
    \label{eq:rebuild}
\end{equation}

\subsubsection{The Shifted-Force Potential}

To ensure both energy and force vanish at $r_c$:
\begin{equation}
    \phi_{\text{SF}}^*(r^*) = \phi^*(r^*) - \phi^*(r_c^*) + \left.\frac{d\phi^*}{dr^*}\right|_{r_c^*}\!\!(r^* - r_c^*)
    \label{eq:shifted_force}
\end{equation}

% [EXPLAIN IN WORDS — what does subtracting the linear term do?]

\subsubsection{NumPy Vectorization}

Instead of nested loops, all pair displacements are computed simultaneously using broadcast array operations, reducing Python-level overhead by orders of magnitude.

\subsection{Test Results}

% Confirm:
% 1. The optimized force routine gives identical energies to the naive version (for small N).
% 2. Show how often the neighbor list is rebuilt (e.g., every ~X steps).
% 3. (Optional) timing comparison: naive loop vs. vectorized.

% [INSERT RESULTS / TIMING TABLE HERE]

\begin{table}[H]
    \centering
    \caption{Comparison of naive and optimized force calculation for $N = 256$ atoms.}
    \label{tab:m4_timing}
    \begin{tabular}{lcc}
        \toprule
        \textbf{Method} & \textbf{Time per step (ms)} & \textbf{Energy match?} \\
        \midrule
        Naive (all pairs) & [value] & --- \\
        Neighbor list + vectorized & [value] & [Yes/No] \\
        Speedup factor & [ratio] & \\
        \bottomrule
    \end{tabular}
\end{table}

% [WRITE DISCUSSION HERE]
